%% notes-tex/028-gene-essentiality/tables/dcell_style.tex
%% GENERATED by experiments/028-gene-essentiality/scripts/notes_tex_tables.py
%% SOURCE: results/dcell_style_*.json
\begin{table}[H]
  \centering
  \footnotesize
  \caption[Round 108: DCell over these splits]{Round 108, DCell's ontology hierarchy on this endpoint. Every term of the Gene Ontology gets a small subsystem whose state is a function of its children's states and its own gene input, wired along the same 2,362-term multi-hot the graph baselines read, with DCell's auxiliary per-term heads at weight 0.3. The published input gives each gene its own weight into each term, which a held-out gene's deletion never trains, so the first two rows replace it with one shared weight on whether the deleted gene carries the term. Columns: seeds on the released split, their mean held-out AUROC and its spread, the five-fold ensemble on the same 195 genes, and the paired difference against graph convolution on the same annotations over ten repeated outer splits with the count of splits the row wins. Test blocks overlap across splits, so no p-value. Width and normalization were chosen on validation over ten configurations. Written by \texttt{notes\_tex\_tables.py} from \texttt{dcell\_style\_baseline.py}.}
  \label{tab:dcell-style}
  \begin{tabular}{@{}lrrrrl@{}}
    \hline
    model & seeds & cell mean & sd & fold ensemble & against the graph, ten splits \\
    \hline
    DCell hierarchy, GO membership on each term & 6 & 0.825 & 0.026 & \textbf{0.848} & $-0.067$ (0.037), 0/10 \\
    the same, the module as an affine on the root & 3 & \textbf{0.846} & 0.017 & \multicolumn{1}{c}{--} & $-0.032$ (0.034), 2/10 \\
    the same, no auxiliary term losses & 6 & 0.799 & 0.037 & 0.844 & \multicolumn{1}{c}{--} \\
    DCell's published per-gene input & 6 & 0.578 & 0.057 & 0.651 & $-0.294$ (0.034), 0/10 \\
    \hline
    graph convolution on the same GO & 6 & 0.905 & 0.005 & 0.918 & 0.868 over ten splits \\
    \hline
  \end{tabular}
\end{table}
